\scanpage{8}
\begin{center}\rule{1.5cm}{.4pt}\\[-.9ex]\hspace{.25cm}\rule{1.25cm}{.4pt}\end{center}
% L1–5
Soit $S$ \deleted{\unsure{d'un|de}}de cardinal $6$ muni d'une partition \inserted{$R$} de type $(3.3)$, $\unsure{\tau|T}$ l'ens de ses deux classes, $S_i$ $(i\in\tau)$ ses deux classes. Soit $\unsure{\mathcal G|g}$ le groupe des autom de $S$, donc
\diagram{\begin{tikzcd}[ampersand replacement=\&,column sep=1em,row sep=.4em]
1\arrow[r]\&{\displaystyle\prod_{i\in\tau}\unsure{\mathfrak S|\mathfrak G}_{S_i}}\arrow[r]\&\mathcal G\arrow[r]\&\mathfrak S_\tau\arrow[r]\arrow[d,equal]\&1\\
\&\&\&{\mathbb Z/2\mathbb Z}
\end{tikzcd}}

% L7–11
$\mathcal G$ d'ordre $72$. Soit $S'\subset\mathcal G$ l'ens des éléments \uline{involutifs} de $\mathcal G$ au dessus de l'élément $\ne1$ de $\mathfrak S_\tau$, i.e. l'ens des scindages de l'ext. précédente.

% L12–17
Pour $i\in\tau$ fixé, il est en corr. $1$-$1$ avec $\deleted{\unsure{\mathrm B|\mathrm I}}\operatorname{Isom}(\mathfrak S_i,\mathfrak S_{\deleted{\unsure{j}}\bar i})$ où $\substack{\bar i\in\tau\\\bar i\ne i}$, ou encore il est en corr. $1$-$1$ avec l'ens des corr. bijectives entre $(\mathfrak S_i)_{i\in\tau}$, i.e. des $\Gamma\subset\prod_{i\in\tau}S_i$ \unsure{telles|tels} que \inserted{$\forall i\in\tau$} $\Gamma\unsure{\mathord{-}|\to} S_i$ soit bij.
% archive sheet number: 7

